\chapter{Vakuumenergie, Inflation und primordialen Fluktuationen}

\section{Vakuumenergie und de-Sitter-Raum}

Im flachen Raum lautet die Friedmann-Gleichung
\begin{equation}
\left(\frac{\dot A}{A}\right)^2=\frac{8\pi G}{3}\rho.
\end{equation}
Für Materie nimmt \(\rho\) wie \(A^{-3}\) ab. Für Strahlung nimmt \(\rho\) wie \(A^{-4}\) ab. Vakuumenergie dagegen verdünnt sich nicht:
\begin{equation}
\rho=\rho_0=\text{constant}.
\end{equation}
Enthält das Universum ausschließlich diese Vakuumenergie, ist die rechte Seite eine Konstante. Bezeichnen wir diese Konstante mit \(H^2\),
\begin{equation}
H^2=\frac{8\pi G}{3}\rho_0,
\qquad
H=\sqrt{\frac{8\pi G}{3}\rho_0}.
\end{equation}
Dann ist die Hubble-Konstante tatsächlich konstant:
\begin{equation}
\frac{\dot A}{A}=H.
\end{equation}
Also gilt
\begin{equation}
\dot A=HA,
\end{equation}
und die Lösung ist exponentiell:
\begin{equation}
A(t)=C e^{Ht}.
\end{equation}
Der Abstand zwischen weit voneinander entfernten, ungebundenen Galaxien wächst exponentiell. Das ist der de-Sitter-Raum.

In diesen Koordinaten lautet die Geometrie
\begin{equation}
d\tau^2=dt^2-e^{2Ht}\left(dx^2+dy^2+dz^2\right).
\end{equation}
Die Vakuumenergie ist hier positiv. Im flachen Raum muss sie positiv sein, denn die linke Seite der Friedmann-Gleichung ist nichtnegativ. Negative Vakuumenergie kann nur auftreten, wenn Krümmungsterme berücksichtigt werden; das ist hier jedoch nicht der betrachtete Fall.

Einstein führte die kosmologische Konstante zu einem anderen Zweck ein: Sie sollte gewöhnliche Materie ausgleichen und das Universum statisch halten.

\section{Das gegenwärtig beschleunigt expandierende Universum}

Der de-Sitter-Raum beschreibt auch die Tendenz des gegenwärtigen Universums zu späten Zeiten. Früher folgte der Skalenfaktor dem für Materie dominierenden Verlauf
\begin{equation}
A(t)\propto t^{2/3}.
\end{equation}
Vor einigen Milliarden Jahren bog die Kurve nach oben ab und begann, schneller anzusteigen. Noch ist der Verlauf nicht exakt exponentiell, weil weiterhin andere Energieformen vorhanden sind; die Entwicklung geht jedoch in diese Richtung.

Beschleunigung bedeutet hier
\begin{equation}
\ddot A>0.
\end{equation}
Dies ist eine beobachtete Eigenschaft der Expansion zu späten Zeiten.

Betrachten wir nun den Horizont. Das Hubble-Gesetz lautet
\begin{equation}
v=HD.
\end{equation}
Bei dem Abstand, bei dem die Rezessionsgeschwindigkeit die Lichtgeschwindigkeit erreicht,
\begin{equation}
D=\frac{c}{H}.
\end{equation}
kann man dies gleichwertig als \(HD=c\) schreiben.\editorialnote{Dies ist eine redaktionelle algebraische Umformulierung der mündlich angegebenen Beziehung \(D=c/H\). Die obere Zeile auf der rechten Tafel in Abbildung~\ref{fig:lecture07-horizon} ist nur teilweise lesbar.}

\begin{figure}[ht]
\centering
\includegraphics[width=0.62\textwidth]{lecture_07_figure_02.png}
\caption{Kosmische Horizontbeziehung \(D=c/H\) bei der de-Sitter-Expansion.}
\label{fig:lecture07-horizon}
\lectureframe{00:30:37}
\end{figure}

Dieser Abstand beträgt etwa zehn bis fünfzehn Milliarden Lichtjahre. Damit liegt er zahlenmäßig nahe am Alter des Universums, wenn man dieses mithilfe von \(c\) in einen Abstand umrechnet.

\begin{classroomqa}
\lecturetimestamp{00:14:10}
\audiencequestion{Ist die ungefähre Gleichheit von Horizontskala und Alter des Universums nicht einfach das, was passiert, wenn die Tangente zufällig durch den Ursprung geht?}
\lecturerresponse{Für \(A(t)\propto t\) wäre das genau richtig. Allgemein gilt es nicht. Im de-Sitter-Raum bleibt \(H\) konstant, während das Alter weiter zunimmt. Die gegenwärtige ungefähre Gleichheit ist also keine mathematische Identität.}
\end{classroomqa}

\section{Was Vakuumenergie bedeutet}

Vakuumenergie wird hier mit den Begriffen der gewöhnlichen Quantenmechanik beschrieben. Für einen harmonischen Oszillator lautet die Energie schematisch
\begin{equation}
E\sim p^2+x^2.
\end{equation}
Klassisch ist der Grundzustand der Gleichgewichtspunkt,
\begin{equation}
x=0,\qquad p=0,\qquad E=0.
\end{equation}
doch die Quantenmechanik erlaubt nicht, dass \(x\) und \(p\) beide exakt auf null festgelegt sind. Die Unschärferelation erzwingt Fluktuationen, sodass die Grundzustandsenergie positiv ist. Für den Oszillator ergibt sich
\begin{equation}
E_0=\frac{1}{2}\hbar\omega.
\end{equation}

Dasselbe Prinzip wird nun auf Felder übertragen. Das Universum ist voller harmonischer Oszillatoren: Schwingungen des elektromagnetischen Feldes und ebenso jedes anderen Feldes. Jedes Feld besitzt Nullpunktsenergie. Für das elektromagnetische Feld lautet die Energiedichte schematisch
\begin{equation}
u_{\mathrm{EM}}\sim E^2+B^2.
\end{equation}
Auch wenn keine realen Photonen vorhanden sind, bleiben irreduzible Nullpunktsfluktuationen bestehen. Das ist eine Erklärung für die Vakuumenergie. Hier kommt es allein darauf an, dass diese Energie vorhanden ist und sich bei der Expansion des Universums nicht verdünnt.

\begin{classroomqa}
\lecturetimestamp{00:16:15}
\audiencequestion{Was ist diese Vakuumenergie, und warum sollte etwas übrig bleiben, wenn Materie und Strahlung entfernt würden?}
\lecturerresponse{Weil das Vakuum der Grundzustand der Quantenfelder ist und ein quantenmechanischer Grundzustand nicht unbedingt die Energie null hat. Harmonische Oszillatoren haben eine positive Grundzustandsenergie, und Felder verhalten sich wie Ansammlungen von Oszillatoren. Selbst wenn keine realen Teilchen vorhanden sind, bleibt also Nullpunktsenergie zurück.}
\end{classroomqa}

Das ist nicht die Hintergrundstrahlung bei drei Grad. Die kosmische Mikrowellenhintergrundstrahlung ist gewöhnliche Strahlung. Ihre Energiedichte nimmt ab, weil die Photonen sich verdünnen und jedes Photon rotverschoben wird. Vakuumenergie ist die verbleibende Nullpunktsenergie, die nach Entfernung der realen Photonen zurückbleibt.

\section{Was \(H^{-1}\) bedeutet}

Kehren wir nun zur Beziehung zwischen dem Kehrwert der Hubble-Konstanten und dem Alter des Universums zurück. Wenn
\begin{equation}
A(t)\propto t,
\end{equation}
dann gilt
\begin{equation}
H=\frac{\dot A}{A}=\frac{1}{t}.
\end{equation}
Allgemeiner gilt, wenn
\begin{equation}
A(t)\propto t^p,
\end{equation}
dann
\begin{equation}
H=\frac{p}{t}.
\end{equation}
Für den materiedominierten Fall \(p=2/3\) gilt
\begin{equation}
H=\frac{2}{3t}.
\end{equation}
In Kosmologien mit Potenzgesetz ist \(H^{-1}\) also von derselben Größenordnung wie das Alter, aber nicht genau gleich dem Alter.

Der de-Sitter-Raum ist anders. Dort gilt
\begin{equation}
A(t)=C e^{\gamma t},
\end{equation}
und damit
\begin{equation}
H=\frac{\dot A}{A}=\gamma,
\end{equation}
eine echte Konstante. Hier ist \(H^{-1}\) nicht das Alter des Universums. Es ist die Zeitkonstante des exponentiellen Wachstums: die Zeit, in der sich der Skalenfaktor mit \(e\) multipliziert. In der mündlichen Erklärung heißt es wiederholt ``verdoppeln'', bevor dies korrigiert wird; präzise ausgedrückt ist es eine Multiplikation mit \(e\). Heute liegt diese Zeitkonstante in der Größenordnung von zehn Milliarden Jahren.

\begin{classroomqa}
\lecturetimestamp{00:23:46}
\audiencequestion{Könnten Sie den Zusammenhang zwischen dem Alter des Universums und dem Kehrwert der Hubble-Konstanten noch einmal erklären?}
\lecturerresponse{Für Potenzgesetze erhält man \(H=p/t\), sodass \(H^{-1}\) bis auf einen Faktor der Größenordnung eins dem Alter folgt. Im de-Sitter-Raum ist \(H\) dagegen konstant, während das Alter weiter zunimmt. Daher ist \(H^{-1}\) nur die exponentielle Zeitkonstante, nicht das Alter.}
\end{classroomqa}

\section{Energie und der kosmische Horizont}

Wenn die Vakuumenergiedichte bei der Expansion des Universums konstant bleibt, liegt die Frage nahe, ob die gesamte Vakuumenergie zunimmt. Die Antwort hier ist eingeschränkt. In der allgemeinen Relativitätstheorie umfasst die Energiebilanz auch einen Gravitationsanteil, weshalb der Begriff der Gesamtenergie schwieriger zu fassen ist als in der elementaren Mechanik. Für die praktische Betrachtung kommt es auf den Horizont an.

Im exakten de-Sitter-Raum hat der kosmische Horizont eine feste Größe,
\begin{equation}
D=\frac{c}{H}
=\frac{c}{\sqrt{8\pi G\rho_0/3}}.
\end{equation}
Das Volumen innerhalb dieses Horizonts ist daher ebenfalls fest. Die Vakuumenergie in diesem Volumen beträgt
\begin{equation}
E_{\mathrm{hor}}=\frac{4\pi}{3}D^3\rho_0.
\end{equation}
\editorialnote{Dies fasst die mündliche Rechnung bei 00:31:06--00:32:23 zusammen: Man nimmt die Horizontgröße, berechnet das Volumen der Kugel und multipliziert mit \(\rho_0\).}
Da sowohl \(D\) als auch \(\rho_0\) konstant sind, ändert sich diese Größe nicht mit der Zeit. Im exakten de-Sitter-Raum bleibt die für einen bestimmten Beobachter jemals sichtbare Vakuumenergie des Universums konstant.

Auch im kosmischen Horizont können Dinge hindurchtreten, ähnlich wie beim Horizont eines Schwarzen Lochs. Sobald sie sich außerhalb befinden, sind sie kausal nicht mehr mit uns verbunden.

\begin{classroomqa}
\lecturetimestamp{00:29:27}
\audiencequestion{Wenn die Vakuumenergiedichte bei der Expansion des Universums konstant bleibt, woher kommt dann die zusätzliche Energie?}
\lecturerresponse{Die Vorlesung beantwortet die Frage im Rahmen der allgemeinen Relativitätstheorie, nicht mit dem Newtonschen Bild eines Kastens. Im exakten de-Sitter-Raum bleibt die Horizontgröße konstant, sodass auch die Vakuumenergie in dem Bereich konstant bleibt, der jemals sichtbar sein kann.}
\end{classroomqa}

\section{Warum die Inflation eingeführt wurde}

Derselbe de-Sitter-Mechanismus wird nun in die ferne Vergangenheit zurückverlegt, allerdings mit einer viel kürzeren Zeitkonstante in der Größenordnung von
\begin{equation}
H^{-1}\sim 10^{-32}\,\mathrm{s}.
\end{equation}
Das ist die Idee der Inflation: Vor allem direkt Beobachtbaren durchlief das Universum eine Phase exponentieller Expansion.

Die grundlegende Motivation sind die Flachheit und Glattheit des Universums. Die räumlichen Schnitte erscheinen sehr flach. Auf großen Skalen ist das Universum außerdem sehr homogen. Noch bemerkenswerter ist, dass die Oberfläche der letzten Streuung bereits glatt war; die Glättung kann also nicht erst während der späteren, sichtbaren Entwicklung stattgefunden haben.

Ein Luftballon, aus dem die Luft abgelassen wurde, ist zerknittert. Bläst man ihn ausreichend auf, wird jeder sichtbare Ausschnitt glatt. Die Analogie mit einer Rosine oder Backpflaume veranschaulicht dasselbe: Ein zerknittertes Objekt erscheint auf jedem sichtbaren Ausschnitt glatt, wenn es enorm gedehnt wird. Genau das soll die Inflation bewirkt haben.

Damit wird auch das Horizontproblem in der hier verwendeten Form beantwortet. Weit voneinander entfernte Bereiche der Oberfläche der letzten Streuung sind für uns beide sichtbar, doch während der späteren, sichtbaren Entwicklung hatte das Licht nicht genug Zeit, um von einem zum anderen zu gelangen. Ihre nahezu einheitliche Dichte und Temperatur legen daher eine Vorgeschichte nahe, in der ein kleinerer Bereich enorm gedehnt wurde.

Die Inflation hat eine zwiespältige Eigenschaft. Sie beseitigt die Abhängigkeit von den Anfangsbedingungen und aus demselben Grund einen großen Teil der Informationen darüber, wie das Universum begonnen hat.

Hier stellt sich die Frage nach der Flachheit in zugespitzter Form erneut. Beobachtungsmäßig scheint der Krümmungsparameter \(k\) null zu sein, und zwar auf etwa ein Prozent genau oder besser. Über den sichtbaren Bereich hinweg erscheint die Geometrie euklidisch. Die Inflation erklärt dies damit, dass der Krümmungsradius viel größer wird als der Bereich, den wir sehen können.

\begin{classroomqa}
\lecturetimestamp{00:40:05}
\audiencequestion{Nehmen wir vor der Inflation ein Gleichgewicht an?}
\lecturerresponse{Nein. Die Idee der Inflation ist fast das Gegenteil. Jede vorher vorhandene Struktur wird so stark gedehnt, dass der sichtbare Bereich die Erinnerung an die früheren Einzelheiten verliert.}
\end{classroomqa}

\section{Ein Skalarfeld und sein Potenzial}

Um zu erklären, wie die effektive Vakuumenergie im frühen Universum viel größer gewesen sein könnte als heute, führen wir ein Skalarfeld ein. Anders als elektrische und magnetische Felder, die Vektoren sind, hat ein Skalarfeld keine Richtung. Es besitzt an jedem Punkt lediglich einen Wert:
\begin{equation}
\phi=\phi(x,y,z,t).
\end{equation}
Dieses Feld ist hypothetisch. Es wird mit keinem bekannten Feld in der Natur identifiziert, und seine Quanten werden mit keinen bekannten Teilchen identifiziert.

Der für uns relevante Anteil seiner Energie ist eine potenzielle Energie,
\begin{equation}
V(\phi).
\end{equation}
Mit ``potenzieller Energie'' ist Energie gemeint, die mit dem Feldwert selbst zusammenhängt, nicht mit der zeitlichen Ableitung des Feldes. Verschiedene Werte von \(\phi\) können verschiedene Energien haben.

\begin{figure}[ht]
\centering
\includegraphics[width=0.68\textwidth]{lecture_07_figure_03.png}
\caption{Qualitatives Skalarpotenzial \(V(\phi)\), das zeigt, dass verschiedenen Feldwerten unterschiedliche potenzielle Energien entsprechen.}
\label{fig:lecture07-generic-potential}
\lectureframe{00:47:08}
\end{figure}

\begin{classroomqa}
\lecturetimestamp{00:46:34}
\audiencequestion{Was stellt die horizontale Achse in der Zeichnung von \(V(\phi)\) dar?}
\lecturerresponse{Die horizontale Achse zeigt den Feldwert \(\phi\), die vertikale Achse die potenzielle Energie \(V\).}
\end{classroomqa}

Ist das Feld räumlich und zeitlich konstant oder im betrachteten Bereich annähernd konstant, dann verhält sich \(V(\phi)\) genau wie Vakuumenergie. Sie verdünnt sich nicht. Daher lassen sich die früheren de-Sitter-Formeln wiederverwenden, indem man \(\rho_0\) durch \(V(\phi)\) ersetzt. Insbesondere gilt
\begin{equation}
H=\sqrt{\frac{8\pi G}{3}V(\phi)}.
\end{equation}

Wenn \(\phi\) von Ort zu Ort variiert, blähen sich verschiedene Bereiche mit unterschiedlichen Zeitkonstanten auf. Ist die Dehnung jedoch groß genug, wird jeder sichtbare Ausschnitt annähernd homogen.

\begin{classroomqa}
\lecturetimestamp{01:05:11}
\audiencequestion{Wirkt \(\phi\) hier wie ein Zeitparameter?}
\lecturerresponse{In dieser Diskussion ja, insofern das Feld beim Herabrollen entlang des Potenzials eine monotone Funktion der Zeit ist.}
\end{classroomqa}

\section{Ein Plateau, eine Klippe und Wiederaufheizung}

Das für diese Erklärung erforderliche Potenzial wird sehr flach, wenn auch nicht vollkommen flach, erreicht dann den Rand einer Klippe, fällt steil ab, erreicht aber nicht ganz null und steigt anschließend wieder an.

\begin{figure}[ht]
\centering
\includegraphics[width=0.68\textwidth]{lecture_07_figure_04.png}
\caption{Inflationspotenzial mit hohem Plateau, steilem Abfall und einem tiefer liegenden Minimum.}
\label{fig:lecture07-inflationary-potential}
\lectureframe{01:02:30}
\end{figure}

Diese Form ist stark konstruiert, wird aber eingeführt, weil die kosmologische Erklärung etwas in dieser Art zu erfordern scheint.

Das Feld beginnt weit oben auf dem Plateau. Da die Steigung gering und die mit der Expansion verbundene effektive Viskosität groß ist, ändert sich \(\phi\) nur sehr langsam. Ein langsam durch Honig fallender Stein trifft die Vorstellung besser als ein Gegenstand, der über einen Tisch gleitet. Je schneller das Universum expandiert, desto größer wird diese Viskosität.

\begin{classroomqa}
\lecturetimestamp{00:59:07}
\audiencequestion{Was bringt \(\phi\) dazu, sich zu bewegen?}
\lecturerresponse{Die Steigung des Potenzials treibt das Feld an. Die Steigung ist gering und die kosmologische Viskosität groß, daher verläuft die Bewegung langsam und nicht schnell.}
\end{classroomqa}

Während \(\phi\) sich langsam über das Plateau bewegt, ändert sich \(V(\phi)\) nur wenig. Deshalb inflationiert das Universum nahezu so, als wäre die Vakuumenergie konstant. Die Zeitkonstante während dieser Phase liegt in der Größenordnung von
\begin{equation}
H^{-1}\sim 10^{-32}\,\mathrm{s}.
\end{equation}
Eine grobe mündliche Schätzung verortet die Plateauhöhe bei etwa \(10^{-14}\) in Planck-Einheiten.\editorialnote{Dies ist die mündliche Schätzung als ``beste Vermutung'' bei 01:14:38--01:14:47 und kein in der Vorlesung hergeleiteter Wert.}

Wie lange dauerte die Inflation? Das ist unbekannt, denn Breite und Flachheit des Plateaus sind nicht bekannt. Bekannt ist, dass sich das Universum mindestens ungefähr fünfzigmal mit \(e\) multiplizierte; sechzig ist die üblicherweise angegebene Zahl.\editorialnote{Redaktionelle Umformulierung der mündlichen Erläuterung zu den E-Faltungen bei 01:14:05--01:19:27.} Als E-Faltungen bezeichnet man diese Faktoren: Jede E-Faltung entspricht einer Multiplikation mit \(e\).

Erreicht das Feld den Rand, ändert sich die potenzielle Energie rasch. Das Feld fällt hinab, seine potenzielle Energie wird zunächst in kinetische Energie und dann in Wärme und Teilchen umgewandelt, und das Universum wird heiß. Die Bezeichnung ``Wiederaufheizung'' ist in dieser Beschreibung etwas irreführend, denn jede zuvor vorhandene Wärmeenergie wäre während der Inflation bereits verdünnt worden.

Das tiefere Minimum liegt nicht bei null. Der dort verbleibende kleine Wert ist mit der gegenwärtigen langsamen exponentiellen Expansion verbunden.

\begin{classroomqa}
\lecturetimestamp{01:19:48}
\audiencequestion{Könnte das Plateau exakt flach sein, sodass Fluktuationen das Feld irgendwie über den Rand stoßen?}
\lecturerresponse{Nein. Es kann nicht vollkommen flach sein. Es muss eine gewisse Neigung haben. Wäre es zu flach, wäre das wiederaufgeheizte Universum zu inhomogen.}
\end{classroomqa}

``Urknall'' wird hier in zwei verwandten Bedeutungen verwendet. Es kann das Wiederaufheizungsereignis selbst bezeichnen. Es kann auch die früheste direkt sichtbare Phase meinen, nämlich die Oberfläche der letzten Streuung. Entscheidend ist, dass alles tatsächlich Beobachtete weit unterhalb des Inflationsplateaus liegt.

\begin{classroomqa}
\lecturetimestamp{01:45:42}
\audiencequestion{Ist die Wiederaufheizung selbst der Urknall?}
\lecturerresponse{Wenn Sie so wollen, ja. Zu diesem Zeitpunkt war das Universum jedoch noch extrem heiß und undurchsichtig. Die Oberfläche der letzten Streuung entsteht viel später, nachdem das Universum so weit abgekühlt ist, dass es durchsichtig wird.}
\end{classroomqa}

\section{Quantenfluktuationen und die Keime der Galaxien}

Die Inflation macht das sichtbare Universum sehr homogen, aber nicht vollkommen homogen. Ohne Quantenfluktuationen würde das Feld völlig glatt werden. Aufgrund der Quantenmechanik variiert es weiterhin geringfügig von Ort zu Ort.

In der zentralen Darstellung ist der gewöhnliche Raum vertikal und der Feldwert \(\phi\) horizontal aufgetragen. Wäre das Feld vollkommen homogen, hätte jeder Punkt im Raum zu einem gegebenen Zeitpunkt denselben Wert \(\phi\), und alle Punkte würden sich gemeinsam auf den Rand zubewegen. Dann würde jeder Bereich im selben Augenblick über die Klippe stürzen. Quantenfluktuationen verhindern diese vollkommen synchrone Bewegung.

Zwei Bilder veranschaulichen den Gedanken. Das eine zeigt eine kleine Erhebung in einem ansonsten homogenen Feld. Das andere eine Reihe von Soldaten, die sich Arm in Arm auf eine Klippe zubewegen. Sind sie perfekt ausgerichtet, stürzen alle gemeinsam hinab. Sind einige etwas voraus und andere etwas zurück, stürzen sie zu unterschiedlichen Zeiten hinunter.

Nun kommt der wichtige Effekt. Ein Bereich, der früher hinabstürzt, wird früher wiederaufgeheizt. Bis ein späterer Bereich hinabstürzt, hatte der frühere bereits Zeit, sich auszudehnen und etwas abzukühlen. Dieser Unterschied führt nach der Wiederaufheizung zu Variationen der Energiedichte.

\begin{classroomqa}
\lecturetimestamp{01:34:00}
\audiencequestion{Sind die später wiederaufgeheizten Bereiche heißer oder kühler als die früheren?}
\lecturerresponse{Die früheren hatten bereits Zeit, sich auszudehnen und abzukühlen. Die späteren sind daher heißer, wenn sie hinabstürzen, und dieser Unterschied trägt zur Dichteinhomogenität bei.}
\end{classroomqa}

Der beobachtete Effekt ist klein:
\begin{equation}
\frac{\delta\rho}{\rho}\sim 10^{-5},
\end{equation}
etwa ein Teil von hunderttausend. Dies sind primordiale Dichtefluktuationen, keine dunkle Energie. Diese Zahl wird hier nicht allein aus der qualitativen Skizze hergeleitet; sie dient dazu, einzuschränken, wie hoch und wie flach das Plateau sein kann.

Sie sind außerdem notwendig. Ohne sie gäbe es keine Galaxien. Nach der Wiederaufheizung verstärkt die Gravitation das anfängliche Muster. Überdichte Bereiche ziehen Materie aus unterdichten Bereichen an. Die überdichten Bereiche werden noch dichter, die unterdichten noch dünner, und daraus entsteht später die Struktur aus Galaxien und Sternen.

Das ist nicht nur eine qualitative Erklärung. Man kann die Quantenfluktuationen eines Feldes in einem expandierenden Universum berechnen, ihre Entwicklung durch die Wiederaufheizung verfolgen und das Ergebnis mit dem aus der Galaxienverteilung rekonstruierten Fluktuationsmuster vergleichen. Die Übereinstimmung ist sehr gut.

Dies führt auch zu einem Punkt über Entropie. Die gravitative Klumpenbildung erscheint kontraintuitiv, verletzt aber den zweiten Hauptsatz nicht. Die Gravitation verhält sich in der statistischen Mechanik ungewöhnlich: Sie begünstigt die Klumpenbildung, und diese erhöht die Entropie, statt sie zu verringern.

\begin{classroomqa}
\lecturetimestamp{01:40:24}
\audiencequestion{Wie schlägt sich die Mikrowellenhintergrundstrahlung in all dem nieder, wenn \(\phi\) kein elektromagnetisches Feld ist?}
\lecturerresponse{Das primordiale Dichtemuster prägt sich der Oberfläche der letzten Streuung auf. Bei der Wiederaufheizung wird die Feldenergie in gewöhnliche Teilchen und Photonen umgewandelt, und die Mikrowellenhintergrundstrahlung trägt später den Abdruck dieser primordialen Fluktuationen.}
\end{classroomqa}

\section{Akustische Wellen und der nächste Schritt}

Zum Schluss wendet sich die Diskussion kurz den akustischen Wellen zu. Sie entstehen vor der späteren gravitativen Klumpenbildung. Die primordialen Fluktuationen entwickeln sich zu akustischen Schwingungen, die am Himmel Verdichtungen bekannter physikalischer Größe hinterlassen.

Das ist wichtig, denn eine bekannte physikalische Größe in bekannter Entfernung liefert einen Maßstab. Misst man den Winkel, unter dem sie auf der Oberfläche der letzten Streuung erscheint, lässt sich anhand der Geometrie feststellen, ob der Raum flach, positiv oder negativ gekrümmt ist. Die Einzelheiten folgen später.

\begin{classroomqa}
\lecturetimestamp{01:47:19}
\audiencequestion{Was ist mit den akustischen Wellen?}
\lecturerresponse{Sie liefern einen Maßstab am Himmel. Ihre bekannte physikalische Größe und der beobachtete Winkel, unter dem sie auf der Oberfläche der letzten Streuung erscheinen, ermöglichen es, die Geometrie des Raums zu bestimmen. Die Einzelheiten werden in der nächsten Vorlesung behandelt.}
\end{classroomqa}